\section[L. In Filtration 3]{\myboxxii{      \BF{\TT{L}}} In Filtration 3}

\textit{Time limit}   2 seconds

\textit{Memory limit} 1 Gigabyte

\textit{Input} \qquad standard input

\textit{Output} \qquad standard output


\begin{quote}
	\textit{除恶务尽——何为「恶」，我们自有裁定。}
\end{quote}

\subsection{Problem Statement}

Maddy is still purging Baddy's evil army\ldots for the third time!

The game takes place on an infinite Cartesian plane, which you may imagine as an infinite chessboard. Baddy controls $n$ stationary \BF{cannons}, placed at $(x_1,y_1),(x_2,y_2),\ldots,(x_n,y_n)$.

Under Chinese chess rules, a cannon attacks by jumping over a screen. If cannons $A$ and $B$ share a row or column and \BF{exactly one} other cannon $C$ lies between them, $A$ attacks $B$ (and $B$ attacks $A$). Cannon $C$ is their ``screen.''

At each step, Maddy may destroy a cannon \BF{that is not currently attacked by any other cannon}. The destroyed cannon is immediately removed; it can no longer attack or serve as a screen.

Maddy wants to destroy all $n$ cannons. Does there exist an order of destruction that clears the board?

\subsection{Input}

The first line contains an integer $T$ ($1 \le T \le 10^5$), the number of test cases.

The first line of each test case contains an integer $n$ ($1 \le n \le 2 \times 10^5$), the number of cannons.

Each of the next $n$ lines contains two integers $x_i,y_i$ ($-10^9 \le x_i,y_i \le 10^9$), the position of cannon $i$.

The sum of $n$ over all test cases does not exceed $2 \times 10^5$. All cannon positions within a test case are distinct.

\subsection{Output}

For each test case, output \TT{YES} if some valid destruction order clears the board, or \TT{NO} otherwise.

The answer is case-insensitive; for example, \TT{yes}, \TT{YES}, \TT{Yes}, and \TT{YeS} are all accepted as affirmative answers.

\subsection{Sample 1}

\begin{minipage}{\linewidth}
	\begin{minipage}[t]{0.49\linewidth}
		\begin{lstlisting}[title=\texttt{standard input}]
3
3
0 0
0 1
0 2
5
0 0
0 1
0 2
1 1
2 1
4
0 0
1 0
2 0
3 0
\end{lstlisting}
	\end{minipage}
	\begin{minipage}[t]{0.49\linewidth}
		\begin{lstlisting}[title=\texttt{standard output}]
YES
YES
NO
\end{lstlisting}
	\end{minipage}
\end{minipage}


\subsection{Sample 1 Explanation}

In the first test case, the three cannons share a column. Initially, $(0,0)$ and $(0,2)$ attack each other over the screen at $(0,1)$, while $(0,1)$ is unattacked. Destroy $(0,1)$ first. Then the other two are unattacked and may be destroyed in either order.

In the third test case, four cannons share a row. Initially, $(0,0)$ and $(2,0)$ attack each other over $(1,0)$; $(1,0)$ and $(3,0)$ attack each other over $(2,0)$. Thus every cannon is attacked, none can be destroyed, and the answer is \TT{NO}.
